\section{Introduction}
\label{sec:intro}

Large language models are increasingly deployed not as conversational chatbots, but as autonomous agents that write code~\citep{swebench}, execute terminal commands, orchestrate tool calls, and engage in multi-turn problem-solving sessions spanning hundreds of interactions. This shift from single-turn chat to multi-step agentic workflows has profound implications for inference infrastructure: a typical coding agent request contains 17,000 input tokens (the repository context) with only 800 output tokens (a code patch), while a chat benchmark assumes 500 input and 300 output tokens---an 11$\times$ difference in input/output ratio.

Despite this transformation, the dominant methodology for evaluating LLM inference systems remains rooted in chat-era assumptions. InferenceX uses random tokens with arbitrary input/output lengths. MLPerf Inference~\citep{mlperf} employs fixed reference datasets designed for hardware comparison rather than workload characterization. Even framework-native benchmarks in vLLM and SGLang~\citep{sglang} primarily use ShareGPT conversation logs, which capture casual chat but not the long-context, prefill-heavy patterns of agentic workloads. As we demonstrate in Section~\ref{sec:results}, this mismatch leads to fundamentally wrong conclusions about system capacity: a server that sustains 920 tokens/s under chat-short workloads may achieve only 240 tokens/s under coding-agent workloads at the same concurrency.

In this work, we make three contributions:

\begin{enumerate}
    \item \textbf{Agentic workload taxonomy and dataset.} We define a four-tier workload taxonomy (Section~\ref{sec:taxonomy}) spanning real agent traces extracted from SWE-Bench and TerminalBench coding sessions, ShareGPT multi-turn conversations, and synthetic stress tests. Each tier captures a distinct region of the input/output length distribution that current benchmarks miss (Figure~\ref{fig:isl_osl}).

    \item \textbf{Saturation analysis across models and hardware.} We benchmark 9 models (dense and MoE, 8B--120B parameters) across NVIDIA A100 and H100 GPUs using vLLM and SGLang, sweeping concurrency from 1 to 320 requests. We show that agentic workloads exhibit earlier saturation, higher TTFT sensitivity, and different optimal tensor parallelism configurations than chat workloads (Section~\ref{sec:results}).

    \item \textbf{Kernel-level roofline analysis.} Using NVIDIA Nsight Compute (ncu) hardware counters, we profile 1,900+ CUDA kernels across Llama-8B and Mixtral-8$\times$7B forward passes, constructing roofline plots that reveal the hardware utilization regime of each kernel class. We show that the throughput gap between agentic and chat workloads is explained by a shift from memory-bound decode kernels to compute-bound prefill GEMMs (Section~\ref{sec:roofline}).
\end{enumerate}

We release our complete dataset---2,400+ benchmark result files, ncu kernel profiles, workload traces, and the open-source benchmarking tool---to enable reproducible evaluation of inference systems under realistic agentic conditions.\footnote{Dataset and code available at the project repository.}
